\section*{Task 1 - RQ1:}

\FloatBarrier
\subsection*{Methodology:}

How do different LLMs perform on cross-domain Text-to-SQL? This research question was approached and answered with a zero-shot benchmark of 11 different configurations of LLMs on the full Spider 1.0 development dataset, with the same prompt, metrics, and statistics being used for each configuration.

We chose these 11 configurations to benchmark by manipulating three different factors: family, size and inference mode. The first family we chose was Qwen 2.5-Coder (Hui et al., 2024), a decoder-only LLM (no explicit reasoning), and within that family we varied the size at the following increments: 0.5B, 1.5B, 3B, and 7B. The second family we chose was Qwen3 (Yang et al., 2025), a reasoning LLM, and within that family we varied the size three times at 1.7B, 4B, and 8B. Qwen3 is a hybrid model, meaning it supports both thinking and non-thinking, and so we were able to alter its inference mode, using the same checkpoints, and without altering the weights (Yang et al., 2025). Thus we obtained three more Qwen3 configurations, one at each aforementioned size, to compare against the adjacent thinking configurations. Additionally we included a 4-bit quantized configuration of Qwen2.5-Coder-7B with AWQ to assess the accuracy and speed cost of weight quantization (Lin et al., 2024). Throughout this report we will be referring to these three factors that discern between and describe the configurations: family, or the training recipe (e.g. Qwen2.5-Coder or Qwen3), size, or the model's parameter count (measured in billions), and specifically for Qwen3, the inference mode, or whether thinking is enabled or disabled for that configuration.

All eight of the checkpoints used have a dense decoder-only transformer architecture. Although Qwen3 is a reasoning LLM, its dense models do share a very similar architecture to Qwen2.5 models (Yang et al., 2025). Therefore the Qwen3 models' major differences from the Qwen2.5-Coder models come from how they were trained, rather than their network architectures. Specifically Coder is Qwen2.5 further trained on about 5.5 trillion tokens of mostly code, making it a code specialized model with no explicit reasoning, while Qwen3 was trained with reinforcement learning for reasoning, merged into a single model that supports both modes and smaller sizes are derived from larger Qwen3 models (Hui et al., 2024; Yang et al., 2025). Since none of the models use a Mixture of Experts (MoE) design, the total and active parameter counts were identical for each checkpoint. Holding the broad architecture relatively constant allowed us to focus comparisons on model family, model size and inference mode. However, these factors were not equally isolated. The most isolated comparison was between the thinking and non-thinking inference modes, as both modes run the exact same Qwen3 checkpoint at each size, making the thinking setting in the chat template the only difference. Comparisons across parameter sizes were also largely isolated, as models in the same family share the same training data and pipeline, with scaling to larger models primarily differing in number and width of their layers. Comparing across families was the least isolated, as these families differ in their training and their training data, and their sizes do not match exactly. Although this comparison targets the effect of training, the combination of these differences means that a performance difference between the families cannot be attributed to any single factor. Instead, from this comparison we can reveal how each training recipe performs as a whole, which still answers if a code specialized model or a more general model trained to reason performs better on a Text-to-SQL task. Furthermore, comparing Qwen2.5-Coder to Qwen3 with thinking disabled, and then against Qwen3 with thinking enabled can help separate the differences in performance due to reasoning itself, and the rest of the differences between the families.

All of the above configurations were run unquantized in bfloat16, our one extension from this was the quantized configuration that used Qwen's official AWQ release of Qwen2.5-Coder-7B-instruct (Qwen Team, 2024). AWQ (Activation aware Weight Quantization) distinguishes roughly the top 1\% of weight channels that matter the most within a model by looking at their size of activations passing through them and protects them from large quantization errors by scaling them up prior to quantization (Lin et al., 2024). After scaling was applied the model weights were then quantized to 4-bit, while computation runs in float16. This extension was done to see if reducing the necessary amount of memory would improve speed without introducing too much loss of accuracy. Therefore, including both the bfloat16 and 4-bit AWQ versions of Qwen2.5-Coder-7B allowed for the assessment of this tradeoff between efficiency and accuracy.

All inference was performed locally. The weights were downloaded from the Hugging Face Hub at pinned revisions in October 2026 and were served with vLLM (Kwon et al., 2023), no API was used. Table 1 shows the exact Hugging Face identifier, revision, parameter count, and precision for each checkpoint, with parameter counts taken from each model's card.

\begin{table}[!htbp]
\centering
\caption{Model Configurations evaluated in Task 1}
\label{tab:1}
\small
\begin{adjustbox}{max width=\linewidth}
\begin{tabular}{@{}lcccccccc@{}}
\toprule
\makecell{Model (Hugging \\ {}Face ID)} & Family & Params & Non-emb. & Modes & Decoding & Max new tok. & Precision & Rev. \\
\midrule
Qwen/Qwen2.5-Coder-0.5B-Instruct & Coder & 0.49B & 0.36B & standard & greedy & 512 & bf16 & ea3f247 \\
Qwen/Qwen2.5-Coder-1.5B-Instruct & Coder & 1.54B & 1.31B & standard & greedy & 512 & bf16 & 2e1fd39 \\
Qwen/Qwen2.5-Coder-3B-Instruct & Coder & 3.09B & 2.77B & standard & greedy & 512 & bf16 & 488639f \\
Qwen/Qwen2.5-Coder-7B-Instruct & Coder & 7.61B & 6.53B & standard & greedy & 512 & bf16 & c03e6d3 \\
Qwen/Qwen2.5-Coder-7B-Instruct-AWQ & Coder & 7.61B & 6.53B & standard & greedy & 512 & 4 bit weights, fp16 & 8e8ed24 \\
Qwen/Qwen3-1.7B & Qwen3 & 1.7B & 1.4B & off / on & greedy / sampled & 512 / 8,192 & bf16 & 70d244c \\
Qwen/Qwen3-4B & Qwen3 & 4.0B & 3.6B & off / on & greedy / sampled & 512 / 8,192 & bf16 & 1cfa9a7 \\
Qwen/Qwen3-8B & Qwen3 & 8.2B & 6.95B & off / on & greedy / sampled & 512 / 8,192 & bf16 & b968826 \\
\bottomrule
\end{tabular}
\end{adjustbox}
\par\vspace{3pt}{\raggedright\footnotesize\textit{Note.} Parameter counts from the model cards (Qwen3 cards round to one decimal); sampled = temperature 0.6, top-p 0.95, top-k 20; thinking seeds 42, 66, 73 (plus 137, 255 for 4B); Rev. = first 7 characters of the Hugging Face commit hash.\par}
\end{table}

All evaluation data was performed on the development set of the Spider 1.0 dataset containing 1,034 natural language questions over 20 databases (Yu et al., 2018). Since this dataset is a cross-domain dataset the development databases do not appear in the training split. Therefore, the models must be able to generalize to database schemas not represented by the Spider training examples rather than answering new questions from databases seen in the training split. Our benchmark is zero-shot, so the training split was not used for either few-shot examples and it also wasn't used for fine tuning or adaptation. Furthermore, we did not use a subsample of the development split, instead every model configuration was evaluated on the whole set of 1,034 development questions. This decision removes the potential for sampling to bias specific models or difficulty levels. The project seed of 42 was fixed prior to running any experiments but had no effect on which questions were being evaluated since the whole split was used. The difficulty of the questions was determined by using Spider's official hardness function script (Yu et al., 2018) and assigned each question a difficulty level of easy, medium, hard, or extra hard depending on the properties of its gold SQL query. The majority of the gold queries in the development set appear twice and are paired with two differently worded questions. The total of 1,034 questions actually corresponds to 551 gold queries. However, these questions that correspond to the same gold query are not fully independent observations and so our statistical analysis must account for this pairing.

For prompting and input we used the same prompt for every configuration and it was fixed before any results were seen. Every model received the same underlying zero-shot prompt as to limit any additional source of variation. The prompt itself was made up of two messages, a system message and a user message. First the system message would say: ``You are an expert in SQLite. Given a database schema and a question, write one SQLite query that answers the question''. The user message would give the database schema under the header ``\#\#\# Database schema'', followed by the question under the header ``\#\#\# Question'', and ended with the single instruction: ``Return only the final SQL query inside a sql code block''. Each database schema was given with the database's CREATE TABLE statements. The statements can then be read directly from each SQLite file, exactly as it was stored, including column table names, column names, column types, and primary and foreign keys. Rajkumar et al. (2022) found that this method of using CREATE TABLE statements was effective in representing relational schemas to LLMs. However, their strongest version provided three example rows for each table. We intentionally did not utilize this exact same format in order to keep the prompts shorter, and give all of the models the same purely schema information, which prevents database contents from becoming an additional source of information. This denial of example values and content does also limit useful information and may reduce overall performance, and therefore is a limitation of the evaluation. The final instruction of the prompt requesting only the final SQL query inside a sql block aided in our extraction of the final SQL from each model's output. For Qwen2.5-Coder and Qwen non-thinking models, the output should be the answer. For Qwen3 thinking models the procedure is different because the output can contain a reasoning trace followed by the final SQL. In this case Qwen3 the output was split at the closing think token, and only the text after that token is searched for SQL, thus all queries during reasoning are ignored. If the think token never occurs we conclude that the entirety of the output token allowance was consumed prior to finishing and it is recorded as having no prediction. After separating the answer from the reasoning trace, the same extraction procedure was applied to all configurations. First the primary extraction searched for the last fenced code block to be taken as prediction. If there was no code block present then a regular expression was used to extract the text from the first SELECT or WITH as in spider the model's predicted SQL query typically begins with one of these keywords. If neither of these methods produced a SQL query, the output was then recorded as having no prediction and scored incorrect. Across all the 19,646 outputs collected from single runs, 19,620 predictions were extracted directly from SQL code blocks, another 10 required the regular expression, and 16 had no prediction. In all 16 cases the reasoning trace reached the maximum output token allowance.

\begin{table}[!htbp]
\centering
\caption{Inference Settings by Mode}
\label{tab:2}
\small
\begin{tabularx}{\linewidth}{@{}>{\raggedright\arraybackslash}p{0.2\linewidth}>{\raggedright\arraybackslash}X>{\raggedright\arraybackslash}X@{}}
\toprule
Setting & \makecell{Coder and Qwen3 \\ {}no thinking} & \makecell{Qwen3 \\ {}thinking} \\
\midrule
Decoding & greedy & sampling: temperature 0.6, top-p 0.95, top-k 20 \\
Max new tokens & 512 & 8,192 \\
Reasoning & off & on \\
Seeds & 42 & 42, 66, 73; plus 137, 255 for 4B; fixed before any thinking run \\
\bottomrule
\end{tabularx}
\par\vspace{3pt}{\raggedright\footnotesize\textit{Note.} Since greedy decoding always returns the same output, the seed 42 does not change anything. Thinking runs sample so their results did vary by seed and were repeated. The context window is the total prompt plus output length vLLM allowed, reduced for 8B because it was too large.\par}
\end{table}

As shown in Table 2 the inference settings were divided into non-thinking and thinking configurations. The non-thinking models used greedy decoding while the thinking models used sampling. The non-thinking models used greedy decoding because the model always selects the highest probability next token. That makes this method of decoding deterministic, so with the same model and prompt, the resulting SQL prediction should always be the same. This is why using the standard seed of 42 doesn't change anything. This allows us to have a deterministic benchmark that can be used for clean size comparisons within families, and more consistent comparisons between families by removing sampling variability from the comparisons. The other settings include a temperature of 0, a top-p of 1, and top-k off with up to 512 new tokens. The temperature being 0 makes the decoding greedy, top-p, and top-k have no impact under greedy decoding, and 512 tokens gives ample space for a SQL query. The thinking models use sampling instead as it allows for the model to select the next token from a pool of high probability options, thus the same question with the same Qwen3 weights, but a different seed could produce a different reasoning trace and potentially a different SQL. The specific settings of temperature = 0.6, top-p = 0.95, and top-k = 20 are directly from the official Qwen3 model card (Qwen Team, 2025). We followed the model card for all settings except for Qwen3 without thinking, and the max token cap for Qwen3 thinking. We changed the Qwen3 without thinking from temperature = 0.7, top-p = 0.8, and top-k = 20, to the greedy decoding settings above to make those runs deterministic and directly comparable to the performance of the Coder runs. We also reduced the max token cap from 32,768 to 8,192 out of necessity. The vLLM's context window had to hold the prompt and the output, and it was set to 16,384 for every model except the 8B, which needed a 10,240 token window, to fit in the available GPU memory. Since the longest prompt was roughly 1,000 tokens, we had to cap it close to 8,000 tokens. This was used at every size so that the amount of time each model was allowed to think did not vary between configurations. However, this cap was only reached in 16 of 11,374 thinking outputs, and between 0 to 3 per run, all of which were scored as incorrect, so any effect works against thinking. For thinking models seeds 42, 66, and 73 were used, and additionally for 4B 137 and 255 were used\footnote{I chose 42 for the initial seed as it is the typical choice coming from ``The Hitchhiker's Guide to the Galaxy''. I then chose seed 66 because of order 66 from Revenge of the Sith.Then I chose 73 because in the Big Bang Theory it is declared to be the best number ever as it is the 21st prime number and its inverse 37 is the 12th prime number, and it is a palindrome in binary. I then chose 137 because in Rick and Morty that is the number of the main dimension or main version of Rick Sanchez that the story follows. Finally I used 255 simply because it is the highest value that can be stored in a single unsigned 8-bite byte in binary.}. For Qwen3 models specifically, reasoning was switched on or off using the model's chat template, as shown in Table 2. We ran these models through Google Colab where we used a single NVIDIA L4 GPU with 24 GB VRAM using vLLM 0.30.0, and the unquantized model configurations loaded in bfloat16 precision. The questions were processed in batches, each of 100 prompts, and vLLM generated every prompt in a batch together. Seconds per question was then calculated by dividing the total inference time of each batch by the number of questions in that batch.

For the evaluation we predominantly used execution accuracy (EX), while also recording exact matching (EM), and component matching (F1). All three of these evaluation metrics were computed with the official test suite SQL evaluation code, the standard Spider evaluator (Zhong et al., 2020). EM is Spider's measurement for comparing the predicted query with the gold query, in order for EM to occur all SQL components must be the same (Yu et al., 2018). F1 also focuses on SQL components, but it breaks the predicted and gold queries into their SQL components and treats each one as a set. Each component is then given an F1 score which reveals which part of a query the model gets incorrect. Like EM this metric only compares structure and cannot show whether a query gets the right answer. EX checks whether the query returns the right answer or not by executing the predicted and the gold query on their database then comparing the returned results. It is not penalized for generating a SQL query that is correct but does not match the structure correctly. Additionally, we calculated EX with the gold query values in the prediction to designate an upper bound for identifying errors caused by incorrect values. While EM and F1 were calculated, recorded and included in some analyses, EX was the primary evaluation metric for this project.

For EX the evaluation code was run on a single Spider database associated with an associated question, and the returned row order was ignored unless the gold query specified some order. Failed outputs were never removed from the evaluation and so the denominator for every configuration was 1,034 questions. If an output did not have a prediction it was scored as incorrect by EX and EM, but if a query timed out or crashed it was scored as incorrect for only EX. If Spider's EM parser was unable to read a prediction, it would be scored as incorrect for EM, but if it still returned the correct result it would be correct for EX. We added two safeguards to the procedure after a few failed attempts. First, a 30 second execution limit was added to prevent queries from running indefinitely within SQLite. This only impacted a single output out of 19,646. Second, EX with gold values tries every way of substitution which can grow exponentially, and so we capped these value combinations at 1,000, or approximately 120 seconds of run time. A total of 28 questions were capped, but all of them kept their ordinary EX result, and for all the unaffected questions our scorer's totals matched that of the official script on all runs.

The main limitation of EX is that on a single database it can mark an incorrect query as correct because it returns the correct result by chance (Zhong et al., 2020; Kim et al., 2025). To defend against this we performed a pessimistic EX sensitivity analysis. This entailed taking 73 questions whose gold result was 0, empty, or NULL, and for these questions if a prediction passed EX but was rejected by EM then its EX was marked as incorrect. This yields the worst case scores, attempting to limit the effect of luck. For error analysis we placed incorrect predictions into mutually exclusive groups in the following order: no prediction, timeout, crash, value-only, or structural.

For each accuracy estimate we used a 95\% Wilson confidence interval to assess uncertainty (Wilson, 1927). For our individual comparisons between configurations we used the exact McNemar test since each configuration was evaluated on the same questions (McNemar, 1947; Dietterich, 1998). Afterwards a Holm correction was applied within each metric to account for performing multiple comparisons (Holm, 1979). We also calculated 95\% bootstrap intervals on each paired accuracy difference using 2,000 resamples with seed 42. Instead of resampling all of the questions, we resampled the 551 distinct gold query groups, so that the two questions sharing a gold query stayed together, and the intervals were not made too narrow (Miller, 2024). Since thinking runs use sampling and have several seeds we report the mean with a standard deviation across seeds, and the 13 planned comparisons were all repeated with pessimistic EX.

\FloatBarrier
\subsection*{Results:}

\begin{table}[!htbp]
\centering
\caption{Summary Statistics}
\label{tab:3}
\small
\begin{adjustbox}{max width=\linewidth}
\begin{tabular}{@{}lccccccccccc@{}}
\toprule
Config & EM easy & EM med & EM hard & EM extra & EM all & EX easy & EX med & EX hard & EX extra & \makecell{EX all \\ {}[95\% CI]} & EX seeds \\
\midrule
Coder 0.5B & 59.3 & 25.3 & 23.6 & 6.0 & 30.1 & 65.3 & 46.0 & 36.8 & 22.3 & 45.3 [42.3, 48.3] &  \\
Coder 1.5B & 56.9 & 30.7 & 18.4 & 7.2 & 31.1 & 73.4 & 62.3 & 35.1 & 33.7 & 55.8 [52.8, 58.8] &  \\
Coder 3B & 71.0 & 26.0 & 29.3 & 4.8 & 33.9 & 83.1 & 72.2 & 59.8 & 42.2 & 67.9 [65.0, 70.7] &  \\
Coder 7B & 78.6 & 49.8 & 43.7 & 12.0 & 49.6 & 91.5 & 85.4 & 75.3 & 53.0 & 80.0 [77.4, 82.3] &  \\
Coder 7B AWQ & 77.0 & 51.8 & 43.7 & 9.0 & 49.6 & 91.9 & 86.1 & 77.0 & 51.2 & 80.4 [77.8, 82.7] &  \\
Qwen3 1.7B off & 54.8 & 11.7 & 9.8 & 1.2 & 20.0 & 83.9 & 63.0 & 44.3 & 33.7 & 60.2 [57.1, 63.1] &  \\
Qwen3 1.7B on & 81.9 & 38.8 & 23.0 & 10.2 & 41.9 & 88.3 & 80.0 & 52.9 & 44.6 & 71.8 [68.9, 74.4] & 72.0 ± 0.3 \\
Qwen3 4B off & 66.1 & 32.7 & 17.8 & 6.6 & 34.0 & 85.9 & 76.9 & 58.0 & 47.6 & 71.2 [68.3, 73.9] &  \\
Qwen3 4B on & 84.3 & 44.4 & 31.0 & 15.7 & 47.1 & 91.9 & 84.1 & 69.5 & 58.4 & 79.4 [76.8, 81.8] & 79.0 ± 0.3 \\
Qwen3 8B off & 62.5 & 28.9 & 10.9 & 2.4 & 29.7 & 85.5 & 80.3 & 62.6 & 51.8 & 74.0 [71.2, 76.6] &  \\
Qwen3 8B on & 85.1 & 42.6 & 27.6 & 9.6 & 45.0 & 89.5 & 83.9 & 71.3 & 57.8 & 78.9 [76.3, 81.3] & 79.3 ± 0.5 \\
\bottomrule
\end{tabular}
\end{adjustbox}
\par\vspace{3pt}{\raggedright\footnotesize\textit{Note.} n per difficulty level; easy = 248, medium = 446, hard = 174, extra hard = 166; EX interval is Wilson 95\%\par}
\end{table}

Table 3 shows EM and EX scores for all 11 model configurations across all four Spider question difficulty levels. All configurations were evaluated on the same 1,034 questions which consisted of 248 easy, 446 medium, 174 hard, and 166 extra hard questions. For the configurations with thinking enabled the results from seed 42 are shown with the mean and standard deviation of EX across seeds in the last column. All configurations show a decrease in EM and EX as the difficulty level of the questions increases, with the exception of EM increasing from medium to hard questions on Coder 3B. The extra hard difficulty questions expectedly proved the most challenging with a maximum of 58.4\% EX from Qwen3-4B thinking on. For the unquantized model configurations the Coder-7B had the highest EX score and EM score (80.0\% and 49.6\%), and its 4-bit quantized variation Coder-7B-AWQ had the highest EX (80.4\%) with the same EM.

\begin{figure}[!htbp]
\centering
\caption{Accuracy by Model Size}
\label{fig:1}
\includegraphics[width=\linewidth]{size_curves.png}
\par\vspace{3pt}{\raggedright\footnotesize\textit{Note.} Thinking points are seed means and the x axis uses each model's real parameter count so Coder and Qwen3 sizes don't line up exactly\par}
\end{figure}

Figure 1 and Table 3 show that overall the EM and EX increase with increasing size of the model for Coder configurations, with an EX of 45.3\% at 0.5B parameters to 80.0\% at 7B parameters. Within the Qwen3 family when comparing identical size configurations, the configurations with thinking enabled have a greater EM and EX for every size. Within the thinking off configurations of the Qwen3 family the EX increases as size increases. For thinking on there was also an increase in EX as size increased across seeds (72.0\% vs 79.0\% vs 79.3\%), but notably there is a slight decrease in EX from 4B to 8B on just seed 42 (79.4\% vs 78.9\%). At similar sizes Qwen3 had a higher EX than Coder configurations at the small and medium sizes (1.5B ≈ 1.7B, and 3B ≈ 4B) slightly without thinking (60.2\% vs 55.8\% and 71.2\% vs 67.9\%), and by more with thinking (71.8\% and 79.4\%), while at the highest size (7B ≈ 8B) Coder was higher than both (80.0\% vs 74.0\% without thinking and 78.9\% with thinking).

\begin{figure}[!htbp]
\centering
\caption{Execution Accuracy by Difficulty}
\label{fig:2}
\includegraphics[width=\linewidth]{accuracy_by_difficulty.png}
\par\vspace{3pt}{\raggedright\footnotesize\textit{Note.} Thinking bars are seed means and so their heights differ slightly from the seed 42 values in Table 3\par}
\end{figure}

Figure 2 visually portrays the decrease in EX from easy to extra hard questions for all configurations. Coder-7B-AWQ, and Coder-7B, the two best configurations on EX and EM, stay within about 2 points of each other at every difficulty level.

\begin{table}[!htbp]
\centering
\caption{Planned Comparisons}
\label{tab:4}
\small
\begin{adjustbox}{max width=\linewidth}
\begin{tabular}{@{}lccccc@{}}
\toprule
Comparison & EX diff [CI] & EX Holm p & \makecell{Pessimistic \\ {}Holm p} & EM diff [CI] & EM Holm p \\
\midrule
Coder 0.5B to 1.5B & +10.5 [6.8, 14.2] & \textless{}0.001 & \textless{}0.001 & +1.1 [−2.4, 4.5] & 1.00 \\
Coder 1.5B to 3B & +12.1 [8.7, 15.5] & \textless{}0.001 & \textless{}0.001 & +2.8 [−0.6, 6.4] & 0.36 \\
Coder 3B to 7B & +12.1 [9.0, 15.2] & \textless{}0.001 & \textless{}0.001 & +15.7 [12.1, 19.3] & \textless{}0.001 \\
Coder 1.5B to 7B & +24.2 [20.6, 28.1] & \textless{}0.001 & \textless{}0.001 & +18.5 [14.6, 22.1] & \textless{}0.001 \\
Thinking off to on, 1.7B & +11.6 [8.6, 14.7] & \textless{}0.001 & \textless{}0.001 & +21.9 [18.4, 25.4] & \textless{}0.001 \\
Thinking off to on, 4B & +8.2 [5.4, 11.4] & \textless{}0.001 & \textless{}0.001 & +13.1 [10.0, 16.2] & \textless{}0.001 \\
Thinking off to on, 8B & +4.9 [2.2, 7.8] & \textless{}0.001 & \textless{}0.001 & +15.3 [12.1, 18.6] & \textless{}0.001 \\
Coder 7B vs Qwen3 4B thinking & −0.6 [−3.6, 2.3] & 1.00 & 1.00 & −2.5 [−6.6, 1.3] & 0.38 \\
Coder 7B bf16 vs AWQ & +0.4 [−0.7, 1.5] & 1.00 & 1.00 & 0.0 [−1.6, 1.6] & 1.00 \\
Qwen3 no thinking 1.7B to 4B & +11.0 [7.5, 14.5] & \textless{}0.001 & \textless{}0.001 & +14.0 [10.6, 17.4] & \textless{}0.001 \\
Qwen3 no thinking 4B to 8B & +2.8 [0.2, 5.4] & 0.097 & 0.220 & −4.4 [−7.3, −1.5] & 0.007 \\
Qwen3 thinking 1.7B to 4B & +7.6 [4.8, 10.5] & \textless{}0.001 & \textless{}0.001 & +5.2 [2.2, 8.1] & 0.001 \\
Qwen3 thinking 4B to 8B & −0.5 [−2.4, 1.3] & 1.00 & 1.00 & −2.1 [−4.7, 0.3] & 0.36 \\
\bottomrule
\end{tabular}
\end{adjustbox}
\par\vspace{3pt}{\raggedright\footnotesize\textit{Note.} Differences are the number of questions only the second configuration got right minus the number of questions only the first configuration got right, divided by the total number of questions (EX difference = (only second - only first) / 1,034) The 95\% CI comes from resampling the 551 gold queries. Holm correction was applied separately to the 13 EX tests and the 13 EM tests. The Pessimistic Holm p repeats the EX Holm p test where every EX pass on a gold answer with an EM fail is counted as wrong.\par}
\end{table}

The EX and EM differences between configurations in Table 4 are assessed for significance with exact McNemar tests, which use the questions that only one or the other configuration being compared got right, followed by Holm correction, and our own pessimistic test. The comparisons within the Coder family show that every increase in parameter size is significant on EX (+10.5, +12.1, +12.1, all Holm p \textless{} 0.001). While on EM only 3B to 7B is significant. Comparing the thinking off to on configurations within the Qwen3 family shows that turning thinking on is significant at every size for both metrics. The EX gains from thinking shrink as the size of the configurations grows (+11.6, +8.2, +4.9), and the EM gains from thinking are larger than that of EX at every size (+21.9, +13.1, +15.3), but do not follow any clear trend. Increasing the size of the configuration in the Qwen3 family from 1.7B to 4B is significant on EX for both modes (+11.0 for non-thinking, +7.6 for thinking). Further increasing from 4B to 8B does not show any statistical significance on EX for both non-thinking and thinking. The Coder-7B vs Qwen3-4B thinking comparison and the Coder-7B vs Coder-7B-AWQ comparison are both not statistically significant on differences in either EX or EM. Under the Holm adjusted p values derived from pessimistic EX, all of the comparisons that were significant for initial EX Holm p remained significant, and the single change was on the Qwen3 non-thinking 4B to 8B comparison. Its EX gain of +2.8 was significant independently (raw p = 0.024) prior to correction, but after correction it was no longer significant (p = 0.097), and under pessimistic EX it shrank to a gain of +2.4 with a Holm p = 0.220.

\begin{table}[!htbp]
\centering
\caption{Efficiency Statistics}
\label{tab:5}
\small
\begin{adjustbox}{max width=\linewidth}
\begin{tabular}{@{}lcccc@{}}
\toprule
Config & s / question & \makecell{Mean output \\ {}tokens} & Median & \makecell{Hit token \\ {}limit} \\
\midrule
Coder 0.5B & 0.015 & 40 & 31 & 8 \\
Coder 1.5B & 0.088 & 230 & 253 & 6 \\
Coder 3B & 0.035 & 38 & 35 & 0 \\
Coder 7B & 0.078 & 34 & 28 & 0 \\
Coder 7B AWQ & 0.041 & 34 & 28 & 0 \\
Qwen3 1.7B off & 0.033 & 40 & 35 & 1 \\
Qwen3 1.7B on & 1.27 & 874 & 571.5 & 2 \\
Qwen3-4B off & 0.055 & 38 & 33 & 0 \\
Qwen3-4B on & 2.11 & 867 & 540 & 0 \\
Qwen3-8B off & 0.092 & 38 & 34 & 0 \\
Qwen3-8B on & 5.46 & 1,045 & 644 & 2 \\
\bottomrule
\end{tabular}
\end{adjustbox}
\par\vspace{3pt}{\raggedright\footnotesize\textit{Note.} The seconds per question is calculated from batch time divided by batch size and measures throughput. Mean output tokens include the reasoning trace for thinking configurations. Reached token limit counts the number of outputs that were cut off at the output limit (512 tokens for non-thinking configurations, and 8,192 tokens for thinking configurations)\par}
\end{table}

Table 5 shows that without thinking outputs are short and fast, ranging from 34 to 40 mean tokens for every configuration except for Coder-1.5B, which had an average of 230 tokens and was slower than larger configurations within the same family. Every configuration without thinking took under 0.1 seconds per question. Coder-1.5B also hit the token limit (512) 6 times which was second only to Coder-0.5B. Comparing the thinking and non-thinking Qwen3 models reveals that thinking costs much more, with an average output tokens between 867 and 1,045, and is much longer with an average seconds per question between 1.27 and 5.46. As the size of the thinking configurations grow, thinking time or seconds per question tend to grow alongside it, while token counts don't seem to expand at the same rate (874, 867, and 1,045 tokens, but 1.27, 2.11, and 5.46 seconds). The Coder-7B-AWQ configuration has the same mean output length to Coder-7B but is about 1.9 times faster than its unquantized counterpart (0.041 seconds vs 0.078 seconds).

\begin{figure}[!htbp]
\centering
\caption{Accuracy-Efficiency Tradeoff}
\label{fig:3}
\includegraphics[width=0.75\linewidth]{cost_tradeoff.png}
\par\vspace{3pt}{\raggedright\footnotesize\textit{Note.} Thinking points are averaged across seeds, and output tokens include the reasoning trace for thinking configurations\par}
\end{figure}

Figure 3 visually shows the relationship between accuracy and cost. Coder-7B and Coder-7B-AWQ have the highest EX values, and they also have the fewest output tokens. Qwen3-4B and 8B with thinking both achieve similar EX values, but it takes more than 20 times the amount of tokens to do so. The mode of thinking or non-thinking causes this drastic increase in mean output length; varying the size of the configuration does not appear to have much of an effect (apart from Coder-1.5B).

\begin{table}[!htbp]
\centering
\caption{Return on Thinking}
\label{tab:6}
\small
\begin{adjustbox}{max width=\linewidth}
\begin{tabular}{@{}lccccc@{}}
\toprule
Size & Easy & Medium & Hard & Extra hard & All \\
\midrule
\multicolumn{6}{@{}l}{\textit{Gain, EX points [95\% CI]}} \\
1.7B & +3.5 [−1.0, 8.0] & +17.3 [13.0, 21.8] & +13.2 [5.9, 20.7] & +8.2 [0.4, 16.1] & +11.9 [9.0, 14.8] \\
4B & +5.2 [1.0, 9.6] & +6.5 [2.8, 10.4] & +12.8 [2.9, 22.0] & +9.6 [0.7, 18.3] & +7.8 [5.0, 10.7] \\
8B & +5.0 [1.1, 9.3] & +3.2 [−0.2, 6.7] & +10.3 [1.9, 18.6] & +6.2 [−1.0, 13.8] & +5.3 [2.7, 8.0] \\
\addlinespace
\multicolumn{6}{@{}l}{\textit{Extra output tokens per question}} \\
1.7B & 474 & 613 & 1,234 & 1,498 & 826 \\
4B & 430 & 655 & 1,264 & 1,479 & 836 \\
8B & 508 & 814 & 1,420 & 1,758 & 994 \\
\addlinespace
\multicolumn{6}{@{}l}{\textit{EX points gained per 1,000 extra tokens}} \\
1.7B & 7.4 & 28.3 & 10.7 & 5.5 & 14.4 \\
4B & 12.2 & 10.0 & 10.1 & 6.5 & 9.3 \\
8B & 9.8 & 3.9 & 7.3 & 3.5 & 5.3 \\
\bottomrule
\end{tabular}
\end{adjustbox}
\par\vspace{3pt}{\raggedright\footnotesize\textit{Note.} The thinking results are averaged across all seeds of a size for a question prior to comparison, and so overall gains will differ slightly from the seed 42 comparisons from Table 4. Intervals resample the 551 gold queries and the per-difficulty intervals are wide because each difficulty only has between 166 and 446 questions.\par}
\end{table}

Thinking clearly does positively impact performance on configurations within the Qwen3 family overall and at every size. The largest impact does appear to be on the smallest configuration and decreases as the parameter size grows. The gain itself seems to be largest on medium and hard questions, specifically on medium questions for 1.7B and more so on hard questions for 4B and 8B. There are three per difficulty gains that have a non-positive lower bound, or an interval that includes zero, 1.7B easy, 8B medium, and 8B extra hard, and so these gains are inconclusive. Thinking adds more output tokens at each increasing difficulty level at every size. Ranging from 474 to 508 extra tokens on easy questions, all the way to between 1,498 and 1,758 extra tokens on extra hard questions. The return of these extra tokens (EX per 1,000 extra tokens) is lowest on the extra hard questions at every size (5.5, 6.5, and 3.5), and decreases overall with an increase in model size (14.4, 9.3, and 5.3).

\begin{table}[!htbp]
\centering
\caption{EM Parser failures}
\label{tab:7}
\small
\begin{adjustbox}{max width=\linewidth}
\begin{tabular}{@{}lcccc@{}}
\toprule
Config & \makecell{Cannot \\ {}be parsed} & \makecell{Of those, \\ {}run fine} & \makecell{Of those, \\ {}EX right} & \makecell{EX-correct answers \\ {}lost to the parser} \\
\midrule
Coder 0.5B & 507 (49.0\%) & 124 & 65 & 13.9\% \\
Coder 1.5B & 516 (49.9\%) & 293 & 158 & 27.4\% \\
Coder 3B & 499 (48.3\%) & 379 & 246 & 35.0\% \\
Coder 7B & 328 (31.7\%) & 291 & 208 & 25.2\% \\
Coder 7B AWQ & 336 (32.5\%) & 305 & 214 & 25.8\% \\
Qwen3 1.7B off & 667 (64.5\%) & 543 & 345 & 55.5\% \\
Qwen3 1.7B on & 418 (40.5\%) & 368 & 250 & 33.7\% \\
Qwen3 4B off & 451 (43.6\%) & 410 & 262 & 35.6\% \\
Qwen3 4B on & 403 (39.0\%) & 395 & 267 & 32.5\% \\
Qwen3 8B off & 577 (55.8\%) & 558 & 388 & 50.7\% \\
Qwen3 8B on & 443 (42.9\%) & 441 & 308 & 37.7\% \\
\bottomrule
\end{tabular}
\end{adjustbox}
\par\vspace{3pt}{\raggedright\footnotesize\textit{Note.} Cannot be parsed column counts predictions that Spider's 2018 SQL parser could not read, and thus EM scores as wrong. Of those run counts the ones that are still executed, and of those EX right counts the number that return the correct output.\par}
\end{table}

Table 7 shows that Spider's old parser could not read between 32\% and 65\% of predictions from all the different configurations. EM then scores these as incorrect despite the fact that many of them are capable of running on the database and returning the correct output. This results in between 13.9\% and 55.5\% of each configuration's EX correct answers being scored as incorrect by EM. Showing why EX is used over EM.

\begin{table}[!htbp]
\centering
\caption{Errors and EX Robustness}
\label{tab:8}
\small
\begin{adjustbox}{max width=\linewidth}
\begin{tabular}{@{}lcccccc@{}}
\toprule
Config & Wrong & Crash \% & Value only \% & Structural \% & \makecell{Trivial gold, EX \\ {}right but EM wrong} & \makecell{Pessimistic \\ {}EX} \\
\midrule
Coder 0.5B & 566 & 68.9 & 2.3 & 28.8 & 35 & 41.9 \\
Coder 1.5B & 457 & 49.7 & 5.9 & 44.4 & 43 & 51.6 \\
Coder 3B & 332 & 37.0 & 7.2 & 55.7 & 46 & 63.4 \\
Coder 7B & 207 & 18.4 & 14.5 & 67.1 & 23 & 77.8 \\
Coder 7B AWQ & 203 & 15.8 & 13.8 & 70.4 & 23 & 78.1 \\
Qwen3 1.7B off & 412 & 31.1 & 5.1 & 63.8 & 56 & 54.7 \\
Qwen3 1.7B on & 292 & 19.9 & 9.6 & 69.5 & 40 & 67.9 \\
Qwen3 4B off & 298 & 14.8 & 7.7 & 77.5 & 39 & 67.4 \\
Qwen3 4B on & 213 & 4.2 & 14.1 & 81.7 & 26 & 76.9 \\
Qwen3 8B off & 269 & 7.8 & 12.6 & 79.6 & 43 & 69.8 \\
Qwen3 8B on & 218 & 0.9 & 16.5 & 81.7 & 28 & 76.2 \\
\bottomrule
\end{tabular}
\end{adjustbox}
\par\vspace{3pt}{\raggedright\footnotesize\textit{Note.} Each wrong answer can only be placed in one of the categories and it is addressed in this order: no prediction, timeout, crash, value only, and structural. However minimal no predictions or timeouts so excluded from the table. Trivial gold counts EX passes on the 73 questions whose gold answer is empty or 0 or NULL, thus EM rejects, and pessimistic EX counts all of these as wrong\par}
\end{table}

Table 8 shows that the percent of crashes reduces drastically with an increase in parameter size, and in the Qwen3 family more than Coder. Conversely, structural errors increase as parameter size increases and, again, more in the Qwen3 family. Value errors tend to remain small across all configs, with their largest values being with the configurations with the strongest performances. The pessimistic EX drops every configuration by between 2.2 and 5.4 points, smallest impact on Coder-7B and Coder-7B-AWQ (tied at 2.2), and largest impact on Qwen3-1.7B without thinking. The rankings of all the different configurations remains unchanged when using the pessimistic EX, thus Coder-7B-AWQ still has the best EX and Coder-0.5B still has the worst.

The best of the selected configurations reach between 79.4\% and 80.4\% EX zero shot on Spider's cross-domain development set, and they achieve that in two different ways: either a large code model (Coder-7B) or Qwen3 with thinking at 4B or larger. Making the best performing configurations Coder-7B-AWQ, Coder-7B, Qwen3-4B thinking, and Qwen3-8B thinking. None of these top four configurations significantly differ from one another: Coder-7B vs Qwen3-4B thinking: −0.6 {[}−3.6, 2.3{]}; Holm = 1.00, AWQ vs Coder-7B: +0.4 {[}−0.7, 1.5{]}; Holm = 1.00, and Qwen3 thinking 4B to 8B −0.5 {[}−2.4, 1.3{]}; Holm = 1.00 (Table 4). These two methods of achieving this accuracy do differ greatly in cost. The Coder configurations use substantially less output tokens per question and are much quicker. Coder-7B uses an average of 34 tokens and 0.078 seconds per question, while 4B thinking uses 867 tokens and 2.11 seconds per question, 25.5 times the amount of tokens and 27 times the time (Table 5, Figure 3). In comparison to Coder-7B, Coder-7B-AWQ has the same number of tokens but is faster, 0.041 seconds per question, or 1.9 times faster. In this setting Coder-7B and especially its quantized counterpart Coder-7B-AWQ give the best balance between accuracy and cost. Thinking cannot buy anything that Coder-7B-AWQ can't do but much cheaper, within the tests of the project.

Thinking can substitute for model size but it gives diminishing returns. For example, 1.7B with thinking (71.8\%) has a similar performance or EX, to 4B without thinking (71.2\%), and 4B with thinking (79.4\%) beats 8B without thinking (74.0\%) by 5.4 points (Table 3). Table 4 directly shows that thinking gains are reduced with an increase in size: +11.6, +8.2, +4.9 EX. The best example for this may be the comparison between 4B thinking (79.0\% ± 0.3) and 8B thinking (79.3 ± 0.5) as there is no gain from doubling the size once thinking is on, and it is not as effective. 8B thinking takes 5.46 s per question vs 2.11 s for 4B, about 2.6x more (Table 5), and its return also drops from 9.3 to 5.3 EX points per 1,000 extra tokens (Table 6). Therefore, with thinking on Qwen3-4B is enough, the additional parameters cost more time without adding any additional accuracy. Without thinking neither family consistently performs better than another. Qwen3-1.7B has better accuracy than Coder-1.5B (60.2\% \textgreater{} 55.8\%), and Qwen3-4B has better accuracy than Coder-3B (71.2\% \textgreater{} 67.9\%), but Coder-7B has higher accuracy than Qwen3-8B (80.0\% \textgreater{} 74.0\%) (Table 3). Qwen3's performance advantage at the smaller sized configurations then comes from thinking itself, not the family. Furthermore, differences in size impact the two families differently. In Coder there is a 12.1 point gain that is statistically significant from 3B to 7B (Holm p \textless{} 0.001), while in Qwen3 without thinking gains only 2.8 points from 4B to 8B and it is not statistically significant after correction (Table 4). Code specific training then appears to make the addition of parameters more useful for SQL generation, even if the step sizes are not identical.

The best configurations fail for different reasons than the worse ones. Crashes make up 68.9\% of Coder-0.5B\textquotesingle s errors but only 18.4\% of Coder-7B\textquotesingle s, 4.2\% of Qwen3-4B thinking\textquotesingle s and 0.9\% of Qwen3-8B thinking\textquotesingle s. While structural errors make up 67.1\% to 81.7\% of the top four configurations\textquotesingle{} errors, and value errors are at a high for these top configurations as well, although they stay more consistent across all configurations. Since structural errors make up most of the errors for the top configurations most of what they get incorrect is the logic of the query while still returning the correct output. And the percentage of value errors is slightly higher because these configurations have the lowest numbers of errors (Coder-7B-AWQ = 203, Coder-7B = 207, Qwen3-4B on = 213, and Qwen3-8B on = 218) so as the number of total errors goes down the number of value errors doesn't go down at the same rate so larger coding models and thinking models of at least 4B do not reduce value errors effectively.

Since pessimistic EX lowers every configuration by 2.2 to 5.4 points but does not change the internal ranking of the configurations, and doesn't change any significance of comparisons between configurations (Tables 4 and 8). The only comparison it moves is Qwen3 4B to 8B without thinking, which was already not significant after correction. This suggests that the 8B's small advantage over the 4B without thinking is partly made up of possibly coincidental or lucky passes, and that our results then do not show that the 8B is reliably better than 4B in this mode (Table 4). Overall, about 80\% EX is reachable zero-shot with small open models, and the cheapest way to get there is a quantized 7B code model rather than a reasoning model.

Some limitations for how this research question was approached include the following. EX can accept wrong queries, so true accuracy may be lower for every configuration. Only using one prompt, every model used one zero shot prompt that showed the schema without any sample rows. Only one benchmark split, results are only on the dev split and may not carry over to other benchmarks.

\FloatBarrier
\clearpage
\section*{Task 2 - RQ2:}

\FloatBarrier
\subsection*{Methodology:}

How do the characteristics of reasoning traces relate to Text-to-SQL performance? This research question was approached and analyzed primarily with the 1,034 traces of Qwen3-4B thinking (seed 42) from the Task 1 reasoning run. This specific run was selected over the Qwen3-8B thinking Task 1 reasoning run since it matches it in performance with 2.6 times less time spent per question (Tables 4 and 5). It is also the only size with 5 seeds since it was selected as the base for Task 3 prior to any Task 2 computation. To ensure that the results were not specific to this individual run we replicated it on 10 additional runs across Qwen3-1.7B, 4B, and 8B (4 more seeds for 4B, and 3 seeds for both 1.7B and 8B), yielding 11,374 traces in total. A trace was counted as correct if its final query was correct under EX, not EM, as Task 1 showed that EM compares query structure, not result, and it has an outdated parser (Table 3 and Table 7). Every trace was included, even the unfinished ones that reached the 8,192 token limit (Table 12).

We defined 5 features for each trace. The first feature is trace length, which measures how much reasoning the model did, or the number of reasoning tokens before the closing think tag, as counted by the vLLM. We measure it as log transformed as the counts are right-skewed. It is relevant because harder queries need more planning of joins, nesting, and schema links, and length can signal that a model got lost (Marjanović et al., 2025; Hassid et al., 2025). The second is ``Alternatively'' per 1,000 words, which is the number of times the whole word ``alternatively'' appears in the lowercase trace divided by the total word count, times 1,000. This matters as it marks the model exploring another formulation, such as a JOIN instead of a subquery, and frequent switches have been linked to incorrect answers in math reasoning (Ling et al., 2026; Wang et al., 2025). Third is ``Wait'' per 1,000 words, and this is computed like ``Alternatively'' but it is the number of times ``wait'' appears in the lowercase trace divided by the total word count, times 1,000. This is important because it marks the model stopping to self check a column name, condition or join key (Ling et al., 2026). Fourth is SQL drafts, the number of uppercase, whole word SELECTs found in the trace, which directly counts how often the model drafted or redrafted the query. This is important because this writing and rewriting can reach the right query and write it into a wrong one. Finally the fifth feature is the answer in trace, which is denoted as 1 if the final query appears exactly in the trace after normalization of both, and 0 if not. This is important as it serves as a surface check on faithfulness, on whether the visible reasoning actually contains the answer the model gave (Chen et al., 2025).

We executed 5 different statistical analyses in order to relate each feature to correctness and difficulty. First we used the Mann-Whitney U test with a rank-biserial effect size in order to relate features and correctness. The features are skewed and often zero, so using a ranked test avoids assuming normality. Second we used the Kruskal-Wallis test over the four levels, alongside Spearman correlation with the difficulty coded from 0 to 3, to relate feature and difficulty, this also fits because both are rank based tests, and Spearman uses ordering for difficulty. Third we used AUROC overall and within each difficulty level. This measures how well each feature separates right from wrong answers, and it is able to do this because it is threshold free, and within each level it removes the confound that hard questions take longer on and are more often wrong on. Fourth we used a logistic regression of correctness on difficulty and each feature, which relates to correctness of features with difficulty held fixed. We also added log length to the model to discover if a feature tells us something beyond length. The standard errors were clustered by gold query because paraphrase pairs share gold queries and are thus not independent of each other (Miller, 2024). Fifth we used Fisher\textquotesingle s exact test to relate to answer in trace and correctness. This fits as the answer in trace is a yes or no feature and the group of traces without the answer could potentially be small, and an exact test remains valid with small counts.

\FloatBarrier
\subsection*{Results:}

\begin{table}[!htbp]
\centering
\caption{Descriptive Statistics of Trace Features for Qwen3-4B thinking (seed 42)}
\label{tab:9}
\small
\begin{adjustbox}{max width=\linewidth}
\begin{tabular}{@{}lcccccc@{}}
\toprule
Feature & Group & Easy & Medium & Hard & Extra & All \\
\midrule
n & Correct / incorrect & 228 / 20 & 375 / 71 & 121 / 53 & 97 / 69 & 821 / 213 \\
Reasoning tokens & correct & 290.5 [231, 398] & 385 [289.5, 570.5] & 798 [543, 1,288] & 928 [586, 1,492] & 435 [289, 729] \\
 & incorrect & 550.5 [309.8, 1,342.5] & 896 [529, 1,411.5] & 1,259 [536, 2,238] & 1,419 [929, 2,509] & 1,082 [575, 1,929] \\
\addlinespace
``Alternatively'' per 1,000 words & correct & 0 [0, 0] & 0 [0, 1.9] & 2.3 [0, 4.0] & 2.6 [0, 3.8] & 0 [0, 2.7] \\
 & incorrect & 2.7 [0, 3.9] & 1.3 [0, 3.4] & 2.6 [1.2, 3.8] & 2.2 [1.2, 3.7] & 2.3 [0, 3.8] \\
\addlinespace
SQL drafts (SELECT count) & correct & 1 [1, 2] & 2 [1, 2] & 4 [2, 8] & 4 [2, 8] & 2 [1, 3] \\
 & incorrect & 2 [2, 4] & 2 [2, 4.5] & 6 [2, 14] & 5 [2, 12] & 4 [2, 9] \\
\addlinespace
``Wait'' per 1,000 words & correct & 6.2 [4.9, 8.9] & 5.7 [4.2, 8.1] & 5.6 [4.0, 7.2] & 5.4 [3.5, 6.8] & 5.8 [4.2, 8] \\
 & incorrect & 4.7 [3.8, 8.5] & 6.2 [4.1, 7.9] & 5.1 [3.8, 7.1] & 5.1 [3.2, 7.8] & 5.2 [3.8, 7.7] \\
\addlinespace
Answer in trace (\%) & correct & 98.2 & 95.2 & 95.0 & 84.5 & 94.8 \\
 & incorrect & 100.0 & 98.6 & 92.5 & 97.1 & 96.7 \\
\bottomrule
\end{tabular}
\end{adjustbox}
\par\vspace{3pt}{\raggedright\footnotesize\textit{Note.} Values are medians with interquartile range (IQR) in brackets, except for Answer in trace which is a percentage of traces that contain the final query. Raw token counts are shown for readability, but the tests are run using their log\par}
\end{table}

Table 9 shows the features median and IQR for correct and incorrect traces at each difficulty level. Of the total of 1,034 traces from seed 42, 821 were correct, and 213 were incorrect, and the number of incorrect traces was the lowest at the easy questions. However, the greatest number of incorrect traces was not from the extra hard questions, but this is due to the medium difficulty having the highest volume of questions, 2.6 times the amount of the hard questions, and 2.7 times the amount of the extra hard questions. Incorrect traces were longer than correct ones at every difficulty level, with a median of 551 tokens against a median of 291 tokens on easy, 896 against 385 on medium, 1,259 against 798 on hard, and 1,419 against 928 on extra hard. Correct traces rarely used ``Alternatively'' with a median of none overall and on the easy and medium questions, while incorrect traces used it at every level, and on the hard and extra hard questions they used it about the same between groups (median of 2.3 for correct vs median of 2.6 for incorrect at hard, and 2.6 for correct vs 2.2 incorrect for extra hard). Furthermore, incorrect traces also included more SQL drafts or SELECT counts with a median of 4 vs 2 for all questions, and both groups drafted more on the hard and extra hard questions (Medians ranging from 4 to 6). ``Wait'' appeared similarly between correct and incorrect traces at every level and was slightly more frequent in correct traces overall 5.8 vs 5.2. The answer in trace percentage was high in both overall with a 94.8\% in correct traces and a 96.7\% in incorrect traces. But, the groups did diverge on specifically the extra hard questions, where the correct traces contained their final query less often (84.5\%) than incorrect traces (97.1\%).

\begin{table}[!htbp]
\centering
\caption{Single Feature Tests Against Correctness and Difficulty}
\label{tab:10}
\small
\begin{adjustbox}{max width=\linewidth}
\begin{tabular}{@{}lccccccccccc@{}}
\toprule
 & \multicolumn{2}{c}{Mann–Whitney} & \multicolumn{2}{c}{Kruskal–Wallis} & \multicolumn{2}{c}{Spearman ρ} & \multicolumn{5}{c}{AUROC, wrong vs right} \\
\cmidrule(lr){2-3}\cmidrule(lr){4-5}\cmidrule(lr){6-7}\cmidrule(lr){8-12}
Feature & p & r\_rb & H & p & Difficulty & Length & Easy & Med. & Hard & Extra & All \\
\midrule
Log reasoning tokens & 1.9 × 10⁻³⁵ & 0.55 & 353 & 3.6 × 10⁻⁷⁶ & 0.58 & 1.00 & 0.75 & 0.79 & 0.61 & 0.65 & 0.78 \\
``Alternatively'' per 1,000 words & 6.0 × 10⁻¹⁴ & 0.30 & 134 & 6.7 × 10⁻²⁹ & 0.33 & 0.60 & 0.69 & 0.65 & 0.52 & 0.51 & 0.65 \\
SQL drafts (SELECT count) & 8.3 × 10⁻¹⁹ & 0.38 & 235 & 1.2 × 10⁻⁵⁰ & 0.45 & 0.73 & 0.75 & 0.69 & 0.57 & 0.56 & 0.69 \\
``Wait'' per 1,000 words & 0.037 & -0.09 & 23 & 4.4 × 10⁻⁵ & -0.15 & -0.12 & 0.38 & 0.51 & 0.45 & 0.50 & 0.45 \\
\bottomrule
\end{tabular}
\end{adjustbox}
\par\vspace{3pt}{\raggedright\footnotesize\textit{Note.} AUROC above 0.5 means higher values go with incorrect traces. Rank-biserial r\_rb = 2 x AUROC - 1, so positive values also mean higher values in incorrect traces. Spearman ρ is computed with difficulty coded 0 to 3\par}
\end{table}

In Table 10 it is clear that trace length, SQL drafts, and ``alternatively'' features differed strongly between correct and incorrect traces. With Mann--Whitney p = 1.9 × 10\^{}-35, 8.3 × 10\^{}-19 and 6.0 × 10\^{}-14, with corresponding rank-biserial effect sizes of 0.55, 0.38 and 0.30. Alternatively ``wait'' was able to reach p = 0.037, but its corresponding effect was negligible (r\_rb = -0.09) and in the opposite direction. Every feature differed across difficulty levels (Kruskal-Wallis), and length increased the most with difficulty (Spearman ρ = 0.58), followed by SQL drafts (Spearman ρ = 0.45), and ``alternatively'' (Spearman ρ = 0.33). Yet again ``wait'' did not succeed and fell slightly (Spearman ρ = -0.15). SQL drafts and ``alternatively'' were both correlated with length (ρ = 0.73 and ρ = 0.60), but ``wait'' once again was not (ρ = -0.12). Length separated correct from incorrect traces best, with an AUROC of 0.78 overall, 0.75 on easy questions, 0.79 on medium questions, 0.61 on hard questions, and 0.65 on extra hard. SQL drafts separated them second best with an overall AUROC of 0.69, then ``alternatively'' with 0.65, these latter two overalls were good, however, their hard and extra hard questions were between 0.51 and 0.57, and finally ``wait'' which was 0.45 overall, the worst over all questions. ``Wait'' overall, medium, hard, and extra hard, as well as hard and extra hard questions for SQL drafts and ``alternatively'' all are close to chance, or the same chance as flipping a coin (0.5).

\begin{table}[!htbp]
\centering
\caption{Logistic regression of correctness on each feature, Odds Ratio [95\% CI]}
\label{tab:11}
\small
\begin{adjustbox}{max width=\linewidth}
\begin{tabular}{@{}lll@{}}
\toprule
Feature & \makecell{Difficulty \\ {}+ feature} & \makecell{Difficulty + log \\ {}length + feature} \\
\midrule
Log reasoning tokens & 0.325 [0.247, 0.428], p = 1.4 × 10⁻¹⁵ &  \\
``Alternatively'' per 1,000 words & 0.852 [0.786, 0.923], p = 8.9 × 10⁻⁵ & 0.968 [0.885, 1.059], p = 0.482 \\
SQL drafts (SELECT count) & 0.953 [0.924, 0.983], p = 0.002 & 1.016 [0.991, 1.042], p = 0.221 \\
``Wait'' per 1,000 words & 1.015 [0.962, 1.070], p = 0.589 & 0.977 [0.926, 1.032], p = 0.404 \\
\bottomrule
\end{tabular}
\end{adjustbox}
\par\vspace{3pt}{\raggedright\footnotesize\textit{Note.} Every model includes difficulty. Odds ratios are per one unit of the feature and for log length one unit is about 2.7 times longer. Values below 1 mean lower odds of a correct answer, and standard errors are clustered by gold query\par}
\end{table}

Table 11 shows the logistic regression with difficulty being held fixed. Length had the strongest association of correctness where each 2.7 times longer trace multiplied the odds of a correct answer by 0.33. ``Alternatively'' had an odds ratio (OR) = 0.85 and SQL drafts OR = 0.95, both were associated with incorrect answers. ``Wait'' was not significant (p = 0.59). After log length was added to the model, none of the three features remained significant. For answer in trace Fisher\textquotesingle s exact test found no relation with correctness (p = 0.285).

\begin{table}[!htbp]
\centering
\caption{Replication Across All 11 Thinking Runs}
\label{tab:12}
\small
\begin{adjustbox}{max width=\linewidth}
\begin{tabular}{@{}lcccccccccc@{}}
\toprule
 &  &  & \multicolumn{2}{c}{Median tokens} & \multicolumn{3}{c}{Log length} & \multicolumn{3}{c}{p beyond length} \\
\cmidrule(lr){4-5}\cmidrule(lr){6-8}\cmidrule(lr){9-11}
Run & EX & Unfinished & Right & Wrong & AUROC & OR & p & Altern. & Drafts & Wait \\
\midrule
4B, seed 42 & 79.4 & 0 & 435 & 1,082 & 0.776 & 0.325 & 1.4 × 10⁻¹⁵ & 0.48 & 0.221 & 0.40 \\
4B, seed 66 & 78.7 & 2 & 425 & 991.5 & 0.765 & 0.366 & 5.5 × 10⁻¹⁴ & 0.54 & 0.085 & 0.82 \\
4B, seed 73 & 79.0 & 0 & 431 & 1,067 & 0.780 & 0.287 & 1.4 × 10⁻¹⁸ & 0.75 & 0.112 & 0.27 \\
4B, seed 137 & 78.9 & 1 & 441.5 & 1,013 & 0.768 & 0.326 & 2.6 × 10⁻¹⁶ & 0.53 & 0.067 & 0.25 \\
4B, seed 255 & 78.7 & 0 & 436.5 & 1,087.5 & 0.767 & 0.306 & 6.6 × 10⁻¹⁷ & 0.46 & 0.220 & 0.91 \\
1.7B, seed 42 & 71.8 & 2 & 439.5 & 903 & 0.743 & 0.425 & 2.3 × 10⁻¹⁰ & 0.46 & 0.076 & 0.08 \\
1.7B, seed 66 & 72.0 & 1 & 426 & 937.5 & 0.737 & 0.434 & 2.5 × 10⁻⁹ & 0.49 & 0.218 & 0.23 \\
1.7B, seed 73 & 72.3 & 3 & 441.5 & 976.5 & 0.747 & 0.430 & 7.1 × 10⁻¹⁰ & 0.20 & 4.8 × 10⁻⁴ & 0.26 \\
8B, seed 42 & 78.9 & 2 & 509 & 1,252 & 0.770 & 0.331 & 2.1 × 10⁻¹⁶ & 0.58 & 0.532 & 0.37 \\
8B, seed 66 & 79.1 & 3 & 514 & 1,266 & 0.777 & 0.302 & 9.9 × 10⁻¹⁷ & 0.67 & 0.324 & 0.35 \\
8B, seed 73 & 79.9 & 2 & 513 & 1,215 & 0.766 & 0.337 & 8.0 × 10⁻¹⁵ & 0.37 & 0.183 & 0.07 \\
\bottomrule
\end{tabular}
\end{adjustbox}
\par\vspace{3pt}{\raggedright\footnotesize\textit{Note.} Unfinished counts are traces that reached the 8,192 token limit. Length, AUROC, OR, and p come from the single feature test and the logistic model with difficulty. The last three columns give each feature's p value when added to difficulty and log length\par}
\end{table}

Table 12 repeats the main results on all 11 thinking runs. In every run the incorrect traces were 2.1 to 2.5 times longer than the correct ones, and length remained strongly associated with correctness (AUROC 0.737 to 0.780, OR 0.287 to 0.434, p \textless= 2.5 x 10\^{}-9). The association was weaker at 1.7B (AUROC 0.737 to 0.747) in comparison to at 4B and 8B (0.765 to 0.780). Beyond length, ``alternatively'' and ``wait'' were never significant, and SQL drafts was only significant for one (1.7B seed 73, p = 4.8x10\^{}-4).

\begin{figure}[!htbp]
\centering
\caption{Trace Features by Difficulty and Correctness for Qwen3-4B Thinking (seed = 42)}
\label{fig:4}
\includegraphics[width=\linewidth]{boxplots_full-qwen3-4b-think.png}
\par\vspace{3pt}{\raggedright\footnotesize\textit{Note.} Blue boxes are correct traces and orange are incorrect. Boxes span the IQR with the median as the middle line. Length is shown on the log scale used in its tests. Answer in trace is not plotted because it is a yes/no feature.\par}
\end{figure}

Figure 4 shows the same pattern from Table 9. The boxes for length and SQL Drafts sit higher for incorrect traces at every difficulty level and both rise with difficulty, and the incorrect boxes are also wider, especially for drafts on hard and extra hard questions. While ``Alternatively'' is higher for incorrect traces, mainly on easy and medium questions, and the ``Wait'' boxes overlap at every difficulty level.

\begin{figure}[!htbp]
\centering
\caption{EX by Trace Length Quartile Within Each Level}
\label{fig:5}
\includegraphics[width=0.7\linewidth]{accuracy_by_length_full-qwen3-4b-think.png}
\par\vspace{3pt}{\raggedright\footnotesize\textit{Note.} Each difficulty level's traces are split into four equal groups by length. Each point represents one group placed at its median length on a log scale.\par}
\end{figure}

Figure 5 shows that within the difficulty levels of easy, medium and extra hard, the longest quarter has the lowest EX, in other words, EX falls as traces get longer. (easy = 80.6\%, medium = 62.5\%, extra = 47.6\%, against 98.4\%, 97.3\%, and 73.8\% respectively for the shortest). The only exception is the hard questions. Here the shortest quarter scores 68.2\%, below the second quarter (86.0\%), before EX falls (to 72.1\% and 52.3\%).

Trace length is the only feature that we tested that showed that it relates to correctness and thus Text-to-SQL performance. Its incorrect traces are 2.5 times longer overall and 1.5 to 2.3 times longer within each level (Table 9). Length also separated correct from incorrect answers within each level (AUROC 0.61 to 0.79), and with a fixed difficulty OR = 0.33, p = 1.4 x 10\^{}-15, and this holds for all replications across all three sizes (Table 12).

Therefore, a long trace signals two things at once. It reflects a hard question, harder questions produce longer traces (ρ = 0.58), and a failed attempt, within the same difficulty a longer trace marks an attempt that is going wrong. Thus not only do longer traces not necessarily help, here, they are detrimental and go with incorrect answers. However the true relationship is not simply that shorter is better. On hard questions the shortest quarter of traces scored a 68.2\% and the second quarter scored a 86.0\%, so in that way too little reasoning can fail as well (Figure 5).

``Alternatively'' and SQL Drafts are only related to correctness here because they grow with length. On their own they yield incorrect answers (OR = 0.85, and OR = 0.95, respectively), which is consistent with Wang et al. (2025) who found that frequent thought switching is associated with incorrect answers. Once length is added neither remains significant (p = 0.48, and p = 0.22), and across 11 runs, only SQL drafts was significant for one run (Table 12). Therefore they do not tell us anything that length does not.

``Wait'' carried no useful signal whatsoever (AUROC 0.45, p = 0.59). Its small unadjusted difference has no effect size (Table 10), it fails on slightly harder questions (ρ = -0.15). Muennighoff et al. (2025) found that forcing a model to write "wait" improves accuracy, however that did not predict correctness here. Answer in trace also showed no relation to correctness (94.8\% of correct and 96.7\% of incorrect traces, p = 0.285). This is only a surface check and does not show that the trace produced the answer (Chen et al., 2025).

Some limitations for the approach to answer this research question include the following. This is correlation only, so the analysis cannot actually deliberately show that long reasoning causes errors. One prompt and one family, all traces come from Qwen3 models under one prompt, other reasoning models may have written traces differently. EX false positive could also carry over into which traces count as correct.

\FloatBarrier
\clearpage
\section*{Task 3 - RQ3:}

\FloatBarrier
\subsection*{Methodology:}

To what extent does execution based self consistency (EBSC) improve Qwen3-4B with thinking? We chose to improve Qwen3-4B with thinking seed 42 due to the reasoning configuration from Task 1. This model tied Coder-7B. Thus we want to investigate if a smaller reasoning model can pull ahead in performance with no additional training but just from implementing a strategy to improve reasoning.

We will be using execution based self consistency as our option A method to improve this model's reasoning. Self consistency works by sampling several answers and then keeping the most common one (Wang et al., 2023). Task 1 showed that most of the 4B errors were logic errors, not crashes (Table 8), and so we are aiming to correct some of them by implementing this vote on execution results rather than SQL text, and this is because different queries are capable of returning the same correct answer (Dong et al., 2023). We expected a small overall gain, concentrated on hard and extra hard questions where the room for improvement is the largest. For each question we did the following steps and abided by them as rules: 1. Took the thinking runs of Qwen3-4B from Task 1 and Task 2 (seeds 42, 66, 73, 137, and 255). 2. Executed each query and dropped any that timed out, crashed or were missing. 3. Grouped the remaining queries by identical result rows. 4. Chose the largest group, breaking ties in favor of the group holding the shortest trace, following Task 2 (Hassid et al., 2025). 5. Submitted the trace's query, keeping seed 42's query if every other sample failed.

We also report a three sample version (seeds 42, 66, 73) that we ran as a trial after fixing the rules for the self consistency. We evaluated EBSC on the same 1,034 questions, with the same metrics and scorer as Task 1, overall and by difficulty, against the original seed 42 run. The primary test, or the 5 samples against the single sample use an exact McNemar test with a bootstrap interval, similar to Task 1. Then the three sample comparison is secondary and Holm corrected.

\FloatBarrier
\subsection*{Results:}

\begin{table}[!htbp]
\centering
\caption{Original Single Sample vs EBSC for Qwen3-4B Thinking (Overall and Difficulty)}
\label{tab:13}
\small
\begin{adjustbox}{max width=\linewidth}
\begin{tabular}{@{}lcccccccccc@{}}
\toprule
 & \multicolumn{5}{c}{EX} &  & EX gain & Fixed / & McNemar & Tokens per \\
\cmidrule(lr){2-6}
Solution & Easy & Med. & Hard & Extra & All & EM all & [95\% CI] & broken & p (Holm) & question \\
\midrule
Single sample (seed 42) & 91.9 & 84.1 & 69.5 & 58.4 & 79.4 & 47.1 &  &  &  & 867 \\
EBSC, N = 3 & 91.9 & 84.5 & 75.3 & 60.8 & 80.9 & 48.6 & +1.5 [0.4, 2.9] & 31 / 15 & 0.026 (0.052) & 2,638 \\
EBSC, N = 5 (primary) & 91.9 & 84.1 & 74.1 & 60.8 & 80.6 & 48.9 & +1.2 [-0.3, 2.6] & 32 / 20 & 0.126 & 4,367 \\
Best of 5 (upper bound) & 94.0 & 89.5 & 81.0 & 68.7 & 85.8 &  &  &  &  &  \\
\bottomrule
\end{tabular}
\end{adjustbox}
\par\vspace{3pt}{\raggedright\footnotesize\textit{Note.} EBSC with N=5 votes over seeds 42, 66, 73, 137, and 255, and N=3 over seeds 42, 66, and 73. Fixed/Broken counts the questions the vote turned correct or incorrect compared with the single sample. Tokens per question are summed over all samples.\par}
\end{table}

Table 13 shows that with 5 samples, EX grew from 79.4\% to 80.6\% (+1.2 points, 95\% CI -0.3 to 2.6). Furthermore, EBSC N=5 fixed 32 questions and broke 20, and the gain was not significant (p = 0.126). However, the gain came exclusively from the hard (+4.6) and extra hard (+2.4) questions, while easy and medium remained unchanged. The three sample EBSC gave a slightly better gain (+1.5 points, 95\% CI 0.4 to 2.9, Holm p = 0.052), adding the additional 2 seeds did not help. EBSC used five times the tokens of a single sample (4,367 vs 867 per question), and the best of five samples would reach 85.8\%, so the vote only recovered 1.2 of the 6.4 points available, for a large cost of computation and a much slower model.

\begin{table}[!htbp]
\centering
\caption{One Question EBSC Fixed and One it Broke}
\label{tab:14}
\small
\begin{tabularx}{\linewidth}{@{}>{\raggedright\arraybackslash}p{0.2\linewidth}>{\raggedright\arraybackslash\ttfamily}X@{}}
\toprule
\multicolumn{2}{@{}p{\linewidth}@{}}{\textbf{Fixed, dev 6 (medium). Show the name and the release year of the song by the youngest singer.}} \\
Gold & \texttt{SELECT song\_name , song\_release\_year FROM singer ORDER BY age LIMIT 1} \\
Seed 42 (incorrect) & \texttt{SELECT Name, Song\_release\_year FROM singer WHERE Age = (SELECT MIN(Age) FROM singer)} \\
Vote, 4 of 5 (correct) & \texttt{SELECT Song\_Name, Song\_release\_year FROM singer WHERE Age = (SELECT MIN(Age) FROM singer)} \\
\midrule
\multicolumn{2}{@{}p{\linewidth}@{}}{\textbf{Broken, dev 111 (medium). What is the accelerate of the car make amc hornet sportabout (sw)?}} \\
Gold & \texttt{SELECT T1.Accelerate FROM CARS\_DATA AS T1 JOIN CAR\_NAMES AS T2 ON T1.Id = T2.MakeId WHERE T2.Make = 'amc hornet sportabout (sw)';} \\
Seed 42 (correct) & \texttt{SELECT cars\_data.Accelerate FROM cars\_data JOIN car\_names ON cars\_data.Id = car\_names.MakeId WHERE car\_names.Make = 'amc hornet sportabout (sw)'} \\
Vote, 3 of 5 (incorrect) & \texttt{SELECT cars\_data.Accelerate FROM car\_names JOIN cars\_data ON car\_names.MakeId = cars\_data.Id WHERE car\_names.Make = 'AMC Hornet Sportabout (SW)'} \\
\bottomrule
\end{tabularx}
\par\vspace{3pt}{\raggedright\footnotesize\textit{Note.} Queries are shown as written. In dev 111 the three voting samples wrote the car name in capitals, which matches no rows, so they returned the same empty result\par}
\end{table}

Table 14 shows one question the vote fixed and one that it broke. In dev 6, the single sample returned the singer's name instead of the song name, and the EBSC had a vote where 4 of the 5 seeds chose the correct column and thus correctly returned the song name. In dev 111, three of the five samples wrote the car's name in capitals, which matches no rows, and so in this case they broke the question where seed 42 got it right.

The expectation from the motivation held in direction and location, but it did not hold up in size. EBSC did improve only hard and extra hard questions, as expected, but its overall gain of 1.2 points was small, not significant (p = 0.126), and not worth the extra compute. At 128 times the cost, EBSC was not significantly different from Coder-7B (p = 0.68), which uses 34 tokens per question to attain a similar EX.
